\section*{Open Questions}

The following five questions arose \emph{during this replication}, grounded
in what we actually observed vs.\ what the paper reports. They are
carried, in machine-readable form with follow-on steps, in
\texttt{report/open\_questions.json}.

\begin{description}

\item[Q1.] \textbf{Where exactly does the $15\%$ gap in $Q_{\text{fact}}$
originate?} Our closed-form estimate using the paper's own per-magic-state
spacetime overhead ($5.35\times 10^{5}$ qubit-rounds at $p_g=10^{-4}$)
yields $Q_{\text{fact}}=5.35\times 10^{6}$ data qubits, but the paper's
Table~I reports $6.30\times 10^{6}$ --- a $17.7\%$ upward adjustment.
The paper attributes overhead to (a) an iterative
``balanced-investment'' optimization over per-level surface-code
distances and (b) rounding at each level, but does not publish a line-item
decomposition of the $\sim 15\%$. Which of the two dominates?

\item[Q2.] \textbf{Does the $2/3$-visibility factor in our raw noise
measurement generalize to more realistic $T$-gate injection noise
models?} Our Monte-Carlo raw error slope fit at $0.98$ (interpreted as
$\log$-slope $=1$ contaminated by a $\times 2/3$ visibility bias from
$Z$-error invisibility in the $\{T,T_\perp\}$ measurement) is
consistent, but the paper's assumption $\epsilon_{\text{in}}=0.4 p_g$
implicitly encodes a specific mapping between physical two-qubit gate
error and injected magic-state error rate (Li 2015). How sensitive is
the Table~I factory footprint to the $0.4$ prefactor being, e.g., $0.3$
or $0.5$ in a realistic ion-trap or superconducting implementation?

\item[Q3.] \textbf{Is the paper's assumption $t_{\text{meas/ff}} = 0.1\,t_{sc}$
architecture-realistic beyond distributed ion traps?} The
$11$-hour figure for 1000-bit Shor at $t_{sc}=10^{-5}$~s hinges on
this ratio. Modern superconducting-qubit demonstrators (e.g., Google
2022, Ref [10] of arXiv:2207.06431) report $t_{\text{meas}}$ and $t_{sc}$
that are closer than $1:10$. If $t_{\text{meas/ff}} \sim t_{sc}$ the
factoring runtime blows up by $\times 10$ to $\sim 4.6$ days. What is
the actual state-of-the-art ratio in 2026 on a superconducting logical
qubit?

\item[Q4.] \textbf{How does the O'Gorman-Campbell $6.3$M-qubit estimate
compare to more recent Litinski-style ``game-of-surface-codes'' resource
estimates (2019)?} The paper predates Litinski by 3 years. Litinski's
approach removes the assumption of a monolithic ``factory'' region and
instead sprinkles distillation into idle spacetime. On the same
1000-bit Shor benchmark at $p_g=10^{-4}$, does Litinski's estimate
match, undercut, or exceed O'Gorman-Campbell's $6.3\times 10^{6}$? A
head-to-head reproduction on the same problem instance would clarify
whether the ``time-optimal'' framing is a lower bound or if
Litinski-style trades further reduce it.

\item[Q5.] \textbf{Does the $35 p^{3}$ cubic law's leading-order
truncation break at the boundary of the paper's operating regime?}
The full 15-to-1 output error is $35 p^{3} + O(p^{4})$; the paper
uses only the leading term. At the paper's edge-case
$p_{\text{in}}=10^{-2}$ (top of Table~I), the next-order correction
could contribute $\gtrsim 10\%$ if the $O(p^{4})$ prefactor is not
small. What is the value of the $p^{4}$ coefficient for 15-to-1, and
at what $p_{\text{in}}$ does using $35p^3$ alone become numerically
misleading for factory design?

\end{description}
